%%%%%%%%%%%%
% Question 3
\item Les données suivantes correspondent à la taille (en cm) et au poids (en kg) de 6 filles de 18 ans (tirées de Morgenthaler, 1997, Exercice 3, p.122) :

\begin{comment} $\sum_{i=1}^n x_i = 1013$, $\sum_{i=1}^n x_i^2 = 171\;311$, 
$\sum_{i=1}^n Y_i = 368$, $\sum_{i=1}^n x_iY_i = 62\;625$
x<-c(170 , 167 , 171 , 181 , 158 , 166 )
y<-c(71 , 58 , 56 , 85 , 45 , 53)
c(somy=sum(y),somy2=sum(y^2),somx=sum(x),somx2=sum(x^2))
  somy  somy2   somx  somx2
   368  23600   1013 171311
anova(lm(y~x))
Response: y
          Df Sum Sq Mean Sq F value  Pr(>F)
x          1 863.99  863.99  20.902 0.01025 *
Residuals  4 165.34   41.34          
\end{comment}
\begin{center}
\begin{tabular}{|l|rrrrrr|c|c|}
\hline
&\multicolumn{6}{|c|}{Observations}&Somme & Somme des carrés\\\hline
Taille $x_i$ & 170 & 167 & 171 & 181 & 158 & 166 &1013&$171\,311$\\\hline
Poids $Y_i$ & 71 & 58 & 56 & 85 & 45 & 53&368&$23\,600$\\\hline
\end{tabular}
\end{center}

\begin{itemize}
\item[a)] Calculez le tableau d'analyse de la variance pour ces données. %Encore une fois, je vous invite à le faire d'abord à la main.
\vspace{0.1 in}
\item[b)] Que concluez-vous quant à la relation entre la taille et le poids des filles de 18 ans, au seuil de 5\%?
\end{itemize}

\prof{\begin{enumerate}
\item On trouve $n = 6$, $\sum_{i=1}^n x_iY_i = 62\,625$.\\
On obtient  que

\begin{align*}
S_{xx} &= 171\,311-\dfrac{1013^2}{6} = 282.83\\
S_{yy} &= 23\,600-\dfrac{368^2}{6}=1029.33\\
S_{xy} &= 62\,625 -\dfrac{(1013)(368)}{6} = 494.33
\end{align*}
 
%Les éléments contenus dans tableau d'analyse de la variance sont $ddl_1=p-1=1$, $ddl_2=n-2$, $ddl_{tot}=n-2+p$, $SC_{reg} = \sum_{i=1}^n(\widehat{Y}_i - \overline{Y})^2$, $SC_{res} = \sum_{i=1}^n(Y_i-\widehat{Y}_i)^2$, $SC_{tot} = \sum_{i=1}^n(Y_i - \overline{Y})^2$, $CM_{reg} = SC_{reg}/ddl_1$, $CM_{res} = SC_{res}/ddl_2$ et finalement $F=CM_{reg}/CM_{res}$. \\

En utilisant le fait que $SC_{tot}=S_{yy}$ et $SC_{reg}=S_{xy}^2/S_{xx}$, on peut construire la table ANOVA suivante:



\begin{table}[h!]
\begin{center}
\begin{tabular}{|l|c|c|c|c|}
  \hline
  Source & ddl & SC & CM  & F \\\hline
  Modèle & 1 & 863.99 & 863.99 & 20.9 \\\hline
  Erreur résiduelle & 4 & 165.34 & 41.34 &  \\\hline\hline
  Totale & 5 & 1029.33 &   &   \\
  \hline
\end{tabular}
\end{center}
\end{table}


\item   La statistique $F$ vaut 20.9. Or $F_{0.05,1,4}=7.71$. Puisque $20.9>7.71$, on rejette $H_0:~$``La taille  n'a pas d'effet'' au seuil de 5\% et on conclut que les donn\'ees d\'emontrent que la taille est associée linéairement au poids. \\

\end{enumerate}
}






%%%%%%%%%%
% Question 3
\item 
Voici une sortie statistique (censurée) de la régression linéaire de la pression systolique en fonction de l\rq{}âge d\rq{}un individu. Utilisez les informations fournies pour répondre aux questions suivantes. 

\begin{verbatim}
Analysis of Variance Table

Response: pression

              Df     Sum Sq    Mean Sq      F value       Pr(>F)
age            1     6394.0     6394.0       ????      7.867e-05 ***
Residuals     ??     8393.4     299.8
\end{verbatim}

\begin{itemize}
\item[a)] Quelle est la taille de l'échantillon?
\item[b)] Donnez une estimation de la variance des erreurs du modèle.
\item[c)] Quelle est la valeur de la statistique $F$?
\item[d)] Expliquez comment la valeur $p$ a été calculée.
\item[e)] Que concluez-vous quant au modèle de régression linéaire, au seuil de 5\%?
\end{itemize}

\prof{\begin{enumerate}
\item Puisque $CM_{res} = \frac{SC_{res}}{n-2}$,
$$n = \frac{SC_{res}}{CM_{res}} + 2 = \frac{8393.4}{299.8} + 2 = 29.99666$$
Donc, il y a $n=30$ observations. \\

\item  $\widehat\sigma^2 = CM_{res} = 299.8$. \\

\item  $F = \dfrac{6394}{299.8} = 21.33$\\

\item La valeur $p$ est la probabilité qu'une variable aléatoire ayant une distribution $F$ avec $1$ et $28$ degrés de liberté soit plus grande que la valeur observée de 21.33. \\

\item  Puisque la valeur  $p$ est très petite, on conclut qu'il y a une relation linéaire significative entre l'âge d'un individu et sa pression systolique.
\end{enumerate}}
